\chapter{How to Use This Template}

So far, Section~\ref{chap:Introduction} to \ref{chap:relatedWork} were following the standard order of sections in theses or research projects. That means that all sections from here on would be the actual content of your work, organized in a number of sections with appropriate titles. This might be just one, or a bit more, depending on what you have done. Consult your supervisor with your proposal.

So, since the main content of this template is how to use the template and how to write a successful research report, we now start providing this content, starting with how to use this template. So:

This section focuses exclusively on how to use this template to the best extent.

\section{Setup and Changes}

\begin{itemize}
  \item \textbf{Set your data.} Use \textsc{configuration.tex} to set your name, supervisor, thesis type, and so on. Just go through that file systematically, it should be self-explanatory.
  \item \textbf{Choose your thesis title page.} This is actually \emph{part} of the last point, but I emphasize it here because many miss it: The configuration file contains an entry/question where you have to choose your preferred title page layout. The two offered ones look very different, so choose the one you like the most. (The pre-compiled PDFs folder also provides a preview of these.) Notably, the title page is only shown in the ``long thesis layout'', which is only recommend for long theses. If you only do a 6 pt research project, or maybe a 12 pt one but you still don't have a significant amount of pages, it is recommended to choose the short/compact thesis layout (see next section on PDF compilation), which does not even offer title pages.
  \item \textbf{Check the declaration.} As mentioned in the source code of the declaration document: it is \emph{your} duty to make sure that the declaration text is up to date. Compare against section 18 of \href{https://policies.anu.edu.au/ppl/document/ANUP_004603}{this ANU webpage}. It's most likely still current, but eventually it will change, so check. Should you notice our template to be outdated, change it for yourself and kindly let the author of this document know so that he can update.
  \item \textbf{Change whatever you want.} Keep in mind that this template is not supposed to restrict what you do or how, but \emph{only} to save you time while making the document look good. \textbf{\emph{This means you can make any changes as you wish}}, you do \emph{not} have to adhere this template. For example, you could:
  \begin{itemize}
    \item change the colors used for links, or whether the linked words should be colored themselves (that's chosen by default) or instead remains black but is instead surrounded by a colored box shown only in the PDF but not the print-out,
    \item change the citation style (e.g., use numbers instead of names),
    \item remove the black boxes that are placed after each definition,
    \item or anything else you could think of.
  \end{itemize}
  Note that in principle, you can even discard this template altogether and choose a standard conference layout, such as those provided by AAAI (google AAAI authorkit) or the ACM (again, google ACM authorkit). This is however only advised for 6 pt projects, and checked with your supervisor. Even if you do switch, follow the next point, i.e., read this template and follow its advice!
  \item \textbf{Follow the advice.} Lastly, read this document! And re-read it once you wrote a bit since you can certainly not remember every piece of advice, especially before you even started writing.
\end{itemize}



\section{Building the PDF}

Most crucially, not that we have two very distinct layouts: a compact one that's recommended for 6 pt projects or 12 pt ones without much produced write-up. For all others, especially Honours theses, the long layout should be chosen. You can trivially achieve this, simply by choosing the respective mainfile:
\begin{itemize}
  \item Compile \textsc{main-long.tex} for Honours theses, 24 pt projects, or 12 pt projects that produced a significant amount of write-up.
  \item Compile \textsc{main-short.tex} for any other research project, especially for all 6 pt projects or 12 pt projects that did not produce a significant amount of write-up.
\end{itemize}

Compilation works as always, just open your desired mainfile and compile it! Either via terminal, via your preferred editor (such as Visual Studio Code with \LaTeX{} plugin, such as ``LaTeX workshop''), or via Overleaf.

If you prefer compilation via terminal/command line, we provided a makefile to speed things up. It's set up to use \emph{latexmk} by default. This very useful commandline tool works on all standard operating systems: Linux, MacOS, and Windows. The installation overhead is minimal, and on Linux and macOS it should even already be installed. Check it out here: \url{https://mg.readthedocs.io/latexmk.html}. It's convenient because it takes the least time to compile a document as \LaTeX{} documents might have to be compiled up to \emph{five} times and latexmk knows the exact amount of compilations required.

Instructions on how to use the makefile:
\begin{itemize}
  \item \textbf{Type ``make''}. Just calling ``make'' without arguments will invoke the build tool \emph{latexmk} in its \emph{online mode}, which automatically updates your PDF every single time you save an updated tex file. (So the terminal you used to invoke make will remain forever idle.) Note that when it encounters a compilation error it often requires a terminal input from you. In this case, fix the \LaTeX{} error(s) and then type X to continue. \textbf{This is the preferred mode! So just type ``make''.}
  \item \textbf{Type ``make mk''}. Calling make with the argument \emph{mk} will also invoke \emph{latexmk}, but without its online mode. So you will have to call it each time you want to have an updated PDF.
  \item \textbf{Type ``make all''}. Calling make with the argument ``all'' will invoke the \LaTeX{} and bibtex compiler manually the maximal amount of times to compile the document. This is more time-intense than simply letting latexmk make its job. Not recommended.
  \item  \textbf{Type ``make quick''}. Calling make with the argument ``quick'' will simply compile the document a single time with the \LaTeX{} compiler. This is quick, but is mostly not enough to show the updated PDF. Not recommended.
  \item \textbf{Type ``make clear''} (or ``make clean''). Calling make with the argument ``clear'' or ``clean'' will delete all temporary files, such as .aux, .log and so on. This is very rarely required. However, sometimes when you produce very wrong code, \LaTeX{} compilation simply fails until you delete all these files from previous compilations. In such a rare scenario, ``make clear'' will help.
\end{itemize}

Note that by default the makefile script compiles the long version. To switch it to the short one, call ``make short'' once, it will store that choice in an invisible .buildtarget file.